\documentclass{article}
\usepackage{amsmath,hyperref}
\begin{document}
\section*{A finite-acceleration time lift for the certified cave example}
\paragraph{Exact scope.}
Keep the same two positive geometric paths and the all-late
arclength lower certificate from PR 874, and the same staged
observation interface supported in the model by PR 875.
The vehicle translates with norm of velocity at most $V$ and
norm of acceleration at most $a$ almost everywhere. It starts
and finishes at rest, and stops at the inspection pose. The
body orientation remains fixed. Waiting for a qualified early
cue costs $\tau$ before starting; give the competing late-only
schedulers no such wait. Physical tracking, actual acceleration
calibration and a bound for deployment latency are not asserted.

\paragraph{Positive timings and universal lower bound.}
For a straight segment of length $\ell\ge V^2/a$, the usual
accelerate-at-$a$, cruise-at-$V$, and brake-to-rest profile takes
$\ell/V+V/a$. Apply it along every certified straight chord.
Velocity is zero at every change of direction, so the
piecewise motion obeys the finite acceleration law. Each
tested positive segment satisfies the length premise.
For any late-only trajectory, consider its outbound and return
parts, starting and ending at rest. Even without tracking
turns, the speed profile is bounded above by
$\min(at,V,a(T-t))$ on a part of duration $T$. Each part covers
at least the distance between the two specified poses, enough
for the trapezoidal lower bound. Integrating gives
$T\ge\ell/V+V/a$ on each part. Combining with the geometric
lower bound gives
\[
 T_{\rm late}\ \ge\ 456.0981695/V+2V/a .
\]
This holds over arbitrary trajectories in the same permitted
domain and the same observation envelope, not only the saved
chords. An equal explicit fixed inspection dwell would add to
both arms.

At $V=1$ m/s, $a=5$ m/s$^2$ and $\tau=1$ s, the early
worst-case bound is $452.474308$ s and the late-only lower bound
is $456.4981695$ s. The certified interval of distinguishing
deadlines is therefore nonempty, with width $4.0238615$ s.
It follows from charging the early wait and stopping at all its
turns, not by treating arclength and time as interchangeable.
The limiting early path has 36 straight segments including
return. In that path regime a sufficient acceleration threshold
for $\tau=1$ at $V=1$ is strictly greater than about $3.141$ m/s$^2$;
there is no extrapolation to the faster, lower-acceleration
vehicles for which this stop-at-turn proof may lose separation.

The stipulated one-second source window and serial execution
can be compared against modeled receipt timings, but are not
a field throughput result. This is a time-lift for a declared
ideal plant, no guarantee of actual vehicle health or motion
tracking. The later compound recovery, independent mission,
and I2 acceptance remain open.
\end{document}
